\chapter{Unit Tests} \label{cap:unit}

\section{Why These Are Unit Tests}

The service tests instantiate the class under test directly and replace collaborators with Mockito mocks. This means that a failure points to business logic in the service or renderer rather than to Spring configuration, database state, HTTP routing, or browser behavior. The tests follow the same pattern throughout: create domain objects, program the mocked collaborator responses, execute one service method, then assert the returned object, thrown exception, or repository interaction.

\section{MeetingService}

The \texttt{MeetingServiceTest} class verifies meeting business rules using mocked repositories. This is the correct level for rules such as invitee normalization, overlap validation, default event duration, response validation, and iCal-token lookup because those rules should be tested independently from the database implementation.

The meeting proposal walkthrough is:

\begin{enumerate}
\item Build an organizer and optional invitees as \texttt{User} objects.
\item Configure \texttt{UserRepository} to return real invitees or report missing users.
\item Configure \texttt{MeetingRepository} to represent either a free calendar or an overlapping calendar.
\item Call \texttt{MeetingService.propose(...)} with a title, description, start time, end time, and raw invitee list.
\item Assert that the organizer is accepted, invitees are pending, duplicate and blank entries are ignored, invalid input is rejected, and saving only happens when the proposal is valid.
\end{enumerate}

\newpage 

The invite response walkthrough is:

\begin{enumerate}
\item Configure \texttt{MeetingParticipantRepository.findByMeetingIdAndUserId(...)} for the authenticated user's invite.
\item Call \texttt{MeetingService.respond(...)} with the requested status.
\item Assert that only \texttt{ACCEPTED} and \texttt{DECLINED} are accepted, that the participant is mutated when present, and that missing invites are rejected.
\end{enumerate}

The discovered-event copy walkthrough is:

\begin{enumerate}
\item Build a \texttt{DiscoveredEvent} as it would come from Ticketmaster, SeatGeek, or AgendaLx.
\item Call \texttt{MeetingService.copyFromDiscovered(...)} for the current user.
\item Assert that a missing end time defaults to two hours, the source and venue are preserved in the description, the user is the only accepted participant, and overlapping imports should be rejected.
\end{enumerate}

Together, these tests cover valid and invalid meeting time windows, unknown invitees, organizer auto-acceptance, duplicate and blank invitee filtering, invite response validation, missing invites, discovered event imports, default event duration, and iCal-token lookup. Repositories are mocked so each test can state exactly which repository calls are expected and whether saving should happen. This makes missing overlap checks visible: the current implementation saves meetings without querying for organizer or invitee overlaps, and saves copied discovered events without checking the user's calendar.

\section{ICalService}

The \texttt{ICalServiceTest} class verifies iCalendar rendering rules. This is a pure rendering service, so tests do not need Spring or a database.

The rendering walkthrough is:

\begin{enumerate}
\item Create one or more \texttt{Meeting} objects with participants and statuses.
\item Assign stable ids with \texttt{ReflectionTestUtils} where the generated UID must be predictable.
\item Render the feed with \texttt{ICalService.render(...)}.
\item Assert the VCALENDAR text directly: escaping, CRLF line endings, UTC date formatting, attendee \texttt{PARTSTAT}, and meeting \texttt{STATUS}.
\end{enumerate}

Together, these tests cover iCalendar text escaping, CRLF line endings, UTC date rendering, tentative versus confirmed meeting status, declined attendee \texttt{PARTSTAT} rendering, and RFC 5545 long line folding. The long line folding test is tagged as a bug-revealing test. When run with \texttt{mvn -Pbug-tests test}, it fails and reveals that \texttt{ICalService} appends long fields directly instead of folding them.

\section{UserService}

The \texttt{UserServiceTest} class verifies registration behavior with a mocked \texttt{UserRepository} and \texttt{PasswordEncoder}. It checks that registration should reject blank username, email, and password values before persisting a user.

The walkthrough is deliberately small: the test calls \texttt{UserService.register("", "", "")}, expects an \texttt{IllegalArgumentException}, and verifies that \texttt{UserRepository.save(...)} is never called. This is a server-side validation check, not a browser check.

This test is also tagged as bug revealing. When run with \texttt{mvn -Pbug-tests test}, it fails because \texttt{UserService.register(...)} only checks duplicate usernames and then saves the user with the encoded password. Browser-side \texttt{required} attributes are not sufficient validation because direct HTTP clients can bypass them.
